\slajdid{02-s027}
\begin{frame}[label=02-s027]{Šta se menja NI/NILI transformacijom?}
% element: 02-s027:e01
\Oznaka{Broj kola i integrisanih čipova}\vspace{2mm}
Transformacijom \alert{ne smanjujemo broj logičkih kola}: ukupan broj ostaje isti. Sada koristimo kola \alert{iste vrste}, pa možemo uštedeti na broju potrebnih čipova.
\par\vspace{4mm}
% element: 02-s027:e02
\Oznaka{Broj ulaza i CMOS realizacija}\vspace{2mm}
\alert{Ukupan broj ulaza takođe ostaje isti}. Kod CMOS logičkih kola broj ulaza direktno određuje broj potrebnih tranzistora, a time utiče i na površinu kola.
\par\vspace{4mm}
% element: 02-s027:e03
\Rezultat{Prednost imaju NI kola zbog manje površine.}
\par\vspace{2mm}
Zašto je to tako, objasnićemo kasnije pri proučavanju CMOS logičkih kola.
\end{frame}
